\chapter{Concentration Cells \& Liquid Junction Potential}
\label{chap:concentration_cells_ljp}

\section{Classification and Principles of Concentration Cells}

In an ordinary chemical galvanic cell (such as the Daniell cell), the electromotive force is generated by an inherent difference in the standard chemical potentials ($\mu^\circ$) of two chemically distinct redox couples ($E^\circ_{\text{cell}} \neq 0$).

In stark contrast, a \textbf{Concentration Cell} is an electrochemical device in which:
\begin{enumerate}[leftmargin=*]
    \item Both half-cells are chemically identical in composition (identical electrode metals and identical electrolyte species).
    \item Consequently, the standard cell potential is \textbf{identically zero}:
    \begin{equation}
        \boxed{E^\circ_{\text{cell}} \equiv 0.0000\text{ V}}
    \end{equation}
    \item The electromotive force ($E_{\text{cell}} > 0$) arises \textbf{solely from a difference in concentration} (or activity $a$) of the electrolyte solutions, or from a difference in the pressure (or activity) of the electrode materials themselves.
\end{enumerate}

Concentration cells are divided into two fundamental classes:
\begin{enumerate}
    \item \textbf{Electrode Concentration Cells:} Electrodes have different concentrations/activities and dip into a common electrolyte solution of uniform concentration.
    \item \textbf{Electrolyte Concentration Cells:} Identical electrodes dip into solutions containing the same electrolyte at different concentrations ($c_1 \neq c_2$).
\end{enumerate}

---

\section{Electrode Concentration Cells}

\subsection{1. Amalgam Concentration Cells}
In an amalgam concentration cell, two electrodes consisting of liquid amalgams of the same metal at different activities ($a_1$ and $a_2$, where $a \propto \text{mole fraction of metal dissolved in mercury}$) dip into a single homogeneous aqueous solution of the metal salt:
\begin{equation}
    \ce{M(Hg, } a_1\ce{) | M^{z+}(aq, } c\ce{) | M(Hg, } a_2\ce{)}
\end{equation}
The electrode processes are:
\begin{align}
    \text{Anode (Oxidation):} \quad & \ce{M(Hg, } a_1\ce{) -> M^{z+}(aq) + z e^-} \\
    \text{Cathode (Reduction):} \quad & \ce{M^{z+}(aq) + z e^- -> M(Hg, } a_2\ce{)} \\
    \hline
    \text{Net Cell Reaction:} \quad & \ce{M(Hg, } a_1\ce{) <=> M(Hg, } a_2\ce{)}
\end{align}
Notice that the electrolyte concentration $c$ cancels out completely from the net cell reaction! The process represents solely the spontaneous transfer of metal from the amalgam of higher activity ($a_1$) to that of lower activity ($a_2$).

Applying the Nernst equation ($E^\circ_{\text{cell}} = 0$):
\begin{equation}
    E_{\text{cell}} = 0 - \frac{RT}{zF} \ln \frac{a_2}{a_1} = \boxed{\frac{RT}{zF} \ln \frac{a_1}{a_2} = \frac{2.303 RT}{zF} \log_{10} \frac{a_1}{a_2}}
\end{equation}
\textbf{Spontaneity Criterion:}
For $E_{\text{cell}} > 0$, we must have:
\begin{equation}
    \boxed{a_1 > a_2}
\end{equation}
The amalgam with the higher concentration of metal must always act as the \textbf{Anode} (oxidation).

\subsection{2. Gas Concentration Cells}
A gas concentration cell consists of two gas electrodes operating at different partial pressures ($P_1$ and $P_2$) immersed in the same electrolyte solution:
\begin{equation}
    \ce{Pt, H2(g, } P_1\ce{) | H+(aq, } c\ce{) | H2(g, } P_2\ce{), Pt}
\end{equation}
The electrode processes are:
\begin{align}
    \text{Anode (Oxidation):} \quad & \ce{H2(g, } P_1\ce{) -> 2H+(aq) + 2e^-} \\
    \text{Cathode (Reduction):} \quad & \ce{2H+(aq) + 2e^- -> H2(g, } P_2\ce{)} \\
    \hline
    \text{Net Cell Reaction:} \quad & \ce{H2(g, } P_1\ce{) <=> H2(g, } P_2\ce{)}
\end{align}
The Nernst equation for $n = 2$ is:
\begin{equation}
    \boxed{E_{\text{cell}} = \frac{RT}{2F} \ln \frac{P_1}{P_2} = \frac{0.0591}{2} \log_{10} \frac{P_1}{P_2}} \quad (\text{at } 298\text{ K})
\end{equation}
\textbf{Spontaneity Criterion:}
The cell generates a positive EMF ($E_{\text{cell}} > 0$) when:
\begin{equation}
    \boxed{P_1 > P_2}
\end{equation}
The gas at higher pressure expands spontaneously to the region of lower pressure.

---

\section{Electrolyte Concentration Cells Without Transference}

In an electrolyte concentration cell without transference, two half-cells contain the same electrolyte at different concentrations ($c_1$ and $c_2$) separated by a salt bridge that prevents direct physical contact and liquid junction potentials:

\subsection{1. Reversible with Respect to Cations}
Consider two silver electrodes dipping into $\ce{AgNO3}$ solutions of concentrations $c_1$ and $c_2$ ($c_2 > c_1$):
\begin{equation}
    \ce{Ag(s) | Ag+(aq, } c_1\ce{)} \parallel \ce{Ag+(aq, } c_2\ce{) | Ag(s)}
\end{equation}
The half-cell reactions are:
\begin{align}
    \text{Anode (Oxidation):} \quad & \ce{Ag(s) -> Ag+(aq, } c_1\ce{) + e^-} \\
    \text{Cathode (Reduction):} \quad & \ce{Ag+(aq, } c_2\ce{) + e^- -> Ag(s)} \\
    \hline
    \text{Net Cell Reaction:} \quad & \ce{Ag+(aq, } c_2\ce{) <=> Ag+(aq, } c_1\ce{)}
\end{align}
The reaction is the spontaneous dilution of silver ions from the more concentrated solution ($c_2$) to the more dilute solution ($c_1$).
Applying the Nernst equation ($E^\circ_{\text{cell}} = 0$):
\begin{equation}
    \boxed{E_{\text{cell}} = \frac{RT}{F} \ln \frac{c_2}{c_1} = 0.0591 \log_{10} \frac{c_2}{c_1}} \quad (\text{at } 298\text{ K})
\end{equation}
For a general $z$-valent cation ($\ce{M^{z+}}$):
\begin{equation}
    \boxed{E_{\text{cell}} = \frac{RT}{zF} \ln \frac{c_2}{c_1} = \frac{0.0591}{z} \log_{10} \frac{c_2}{c_1}}
\end{equation}
\textbf{Rule:} For cation-reversible electrodes, the \textbf{more dilute} solution acts as the \textbf{Anode} ($c_1$), and the \textbf{more concentrated} solution acts as the \textbf{Cathode} ($c_2$).

\subsection{2. Reversible with Respect to Anions}
Consider two chlorine gas electrodes dipping into $\ce{KCl}$ solutions of concentrations $c_1$ and $c_2$:
\begin{equation}
    \ce{Pt, Cl2(g) | Cl^-(aq, } c_1\ce{)} \parallel \ce{Cl^-(aq, } c_2\ce{) | Cl2(g), Pt}
\end{equation}
The half-cell reactions are:
\begin{align}
    \text{Anode (Oxidation):} \quad & \ce{2Cl^-(aq, } c_1\ce{) -> Cl2(g) + 2e^-} \\
    \text{Cathode (Reduction):} \quad & \ce{Cl2(g) + 2e^- -> 2Cl^-(aq, } c_2\ce{)} \\
    \hline
    \text{Net Cell Reaction:} \quad & \ce{2Cl^-(aq, } c_1\ce{) <=> 2Cl^-(aq, } c_2\ce{)}
\end{align}
The Nernst equation is:
\begin{equation}
    \boxed{E_{\text{cell}} = \frac{RT}{F} \ln \frac{c_1}{c_2} = 0.0591 \log_{10} \frac{c_1}{c_2}} \quad (\text{at } 298\text{ K})
\end{equation}
\textbf{Rule:} For anion-reversible electrodes, the \textbf{more concentrated} solution acts as the \textbf{Anode} ($c_1$), and the \textbf{more dilute} solution acts as the \textbf{Cathode} ($c_2$).

---

\section{Electrolyte Concentration Cells With Transference \& Liquid Junction Potential ($E_{lj}$)}

When two solutions of the same electrolyte at different concentrations ($a_1$ and $a_2$) are in direct physical contact across a porous barrier without a salt bridge:
\begin{equation}
    \ce{Ag | AgNO3}(a_1) \ \vdots \ \ce{AgNO3}(a_2) | \ce{Ag} \qquad (a_2 > a_1)
\end{equation}
The passage of $1\text{ Faraday}$ of electricity involves not only the electrode reactions, but also the physical migration of ions across the boundary interface (transference).

\begin{figure}[htbp]
\centering
\begin{tikzpicture}[scale=1.1]
    % Boundary container
    \draw[thick] (0,0) rectangle (7,3.5);
    \fill[blue!8] (0.05,0.05) rectangle (3.45,2.7);
    \fill[blue!20] (3.55,0.05) rectangle (6.95,2.7);
    
    % Porous partition
    \draw[line width=2pt, dashed, primarynavy] (3.5,0.05) -- (3.5,2.7);
    \node[above] at (3.5,2.8) {\small \textbf{Porous Boundary}};

    % Electrodes
    \draw[fill=gray!40, thick] (0.8,0.5) rectangle (1.2,3.8);
    \node[above] at (1.0,3.9) {\small $\ce{Ag}$ Anode};
    \draw[fill=gray!40, thick] (5.8,0.5) rectangle (6.2,3.8);
    \node[above] at (6.0,3.9) {\small $\ce{Ag}$ Cathode};

    % Solution labels
    \node at (2.0,1.4) {$\ce{AgNO3}(a_1)$};
    \node at (5.0,1.4) {$\ce{AgNO3}(a_2)$};
    \node at (3.5,-0.3) {$(a_2 > a_1)$};

    % Transference arrows
    \draw[->, line width=1.5pt, cengagegreen] (3.0,0.8) -- (4.2,0.8) node[midway, below] {\footnotesize $t_+ \ce{Ag+}$};
    \draw[->, line width=1.5pt, crimsonred] (4.2,2.0) -- (3.0,2.0) node[midway, above] {\footnotesize $t_- \ce{NO3-}$};
\end{tikzpicture}
\caption{Transference of ions across the liquid-liquid junction in a concentration cell with transference.}
\label{fig:concentration_cell_transference}
\end{figure}

\subsection{Rigorous Derivation of EMF With Transference ($E_{wt}$)}
When $1\text{ F}$ of electric charge passes through the cell:
\begin{enumerate}[leftmargin=*]
    \item \textbf{At the Anode:} $1\text{ mole}$ of $\ce{Ag}$ is oxidized:
    \begin{equation}
        \text{Gain around anode: } +1\text{ mole of } \ce{Ag+}
    \end{equation}
    \item \textbf{Across the Junction:}
    \begin{itemize}
        \item $t_+$ moles of $\ce{Ag+}$ cations migrate from the anode compartment ($a_1$) to the cathode compartment ($a_2$).
        \item $t_-$ moles of $\ce{NO3-}$ anions migrate from the cathode compartment ($a_2$) to the anode compartment ($a_1$).
    \end{itemize}
    \item \textbf{Net Change in Anode Compartment ($a_1$):}
    \begin{equation}
        \Delta n_{\ce{Ag+}} = 1 - t_+ = t_- \text{ moles of } \ce{Ag+} \quad (\text{since } t_+ + t_- = 1)
    \end{equation}
    Along with $t_-$ moles of $\ce{NO3-}$ arriving, exactly $t_-$ moles of $\ce{AgNO3}$ are gained in compartment $a_1$.
    \item \textbf{Net Change in Cathode Compartment ($a_2$):}
    By symmetry, exactly $t_-$ moles of $\ce{AgNO3}$ are lost from compartment $a_2$.
\end{enumerate}
Therefore, the net cell process per Faraday of charge passed is:
\begin{equation}
    t_- \ce{AgNO3}(a_2) \ce{<=>} t_- \ce{AgNO3}(a_1)
\end{equation}
The free energy change is:
\begin{equation}
    \Delta G = t_- \Delta G_{\text{dilution}} = -t_- \cdot 2RT \ln \frac{a_2}{a_1}
\end{equation}
Since $\Delta G = -1 \cdot F E_{\text{wt}}$:
\begin{equation}
    \boxed{E_{\text{wt}} = 2 t_- \frac{RT}{F} \ln \frac{a_2}{a_1}}
\end{equation}
where $t_-$ is the transport number of the anion.

\subsection{Derivation of the Liquid Junction Potential ($E_{lj}$)}
The total potential of the cell with transference ($E_{\text{wt}}$) is the algebraic sum of the electrode potentials without transference ($E_{\text{wot}}$) and the junction potential ($E_{lj}$):
\begin{equation}
    E_{\text{wt}} = E_{\text{wot}} + E_{lj} \implies \boxed{E_{lj} = E_{\text{wt}} - E_{\text{wot}}}
\end{equation}
Recall that for the cell without transference:
\begin{equation}
    E_{\text{wot}} = \frac{RT}{F} \ln \frac{a_2}{a_1}
\end{equation}
Subtracting $E_{\text{wot}}$ from $E_{\text{wt}}$:
\begin{equation}
    E_{lj} = 2 t_- \frac{RT}{F} \ln \frac{a_2}{a_1} - \frac{RT}{F} \ln \frac{a_2}{a_1} = (2t_- - 1) \frac{RT}{F} \ln \frac{a_2}{a_1}
\end{equation}
Since $1 = t_+ + t_- \implies 2t_- - 1 = t_- - t_+$:
\begin{equation}
    \boxed{E_{lj} = (t_- - t_+) \frac{RT}{F} \ln \frac{a_2}{a_1} = (2t_- - 1) \frac{RT}{F} \ln \frac{a_2}{a_1}}
\end{equation}

\begin{conceptbox}[Fundamental Insights into the Liquid Junction Potential]
\begin{itemize}[leftmargin=*]
    \item \textbf{Physical Origin:} $E_{lj}$ arises because the cation and anion diffuse across the liquid boundary at different rates ($u_+ \neq u_-$). The faster-moving ion diffuses ahead, setting up a charge separation and potential gradient across the boundary.
    \item \textbf{When $t_- > t_+$ (e.g., $\ce{HCl}$ with fast $\ce{H+}$? No, for $\ce{HCl}$, $u_+ > u_- \implies t_+ > t_-$):}
    For $\ce{HCl}$, $t_{\ce{H+}} \approx 0.83, t_{\ce{Cl-}} \approx 0.17$. Because $\ce{H+}$ diffuses much faster than $\ce{Cl-}$, the dilute solution becomes positively charged and the concentrated solution negatively charged ($E_{lj} < 0$).
    \item \textbf{When $t_+ = t_- = 0.50$ (The Salt Bridge Principle):}
    \begin{equation*}
        E_{lj} = (0.50 - 0.50) \frac{RT}{F} \ln \frac{a_2}{a_1} \equiv 0
    \end{equation*}
    Because $\ce{K+}$ and $\ce{Cl-}$ have virtually identical mobilities ($t_{\ce{K+}} \approx 0.49, t_{\ce{Cl-}} \approx 0.51$), the diffusion potentials cancel out, completely eliminating $E_{lj}$.
\end{itemize}
\end{conceptbox}

---

\section{Comprehensive Master Solved Illustrations}

\begin{examplebox}[Calculation of Transference EMF and Liquid Junction Potential]
At $25^\circ\text{C}$ ($298\text{ K}$), the EMF of the concentration cell without transference:
\begin{equation*}
    \ce{Ag | AgNO3}(0.01\text{ M}) \parallel \ce{AgNO3}(0.10\text{ M}) | \ce{Ag}
\end{equation*}
is found to be $0.0591\text{ V}$. When the salt bridge is removed and the two solutions are placed in direct contact across a porous partition:
\begin{equation*}
    \ce{Ag | AgNO3}(0.01\text{ M}) \ \vdots \ \ce{AgNO3}(0.10\text{ M}) | \ce{Ag}
\end{equation*}
the measured EMF of the cell with transference is $0.0632\text{ V}$.
Calculate:
\begin{enumerate}[label=(\alph*)]
    \item The transport number of the nitrate ion ($t_{\ce{NO3-}}$).
    \item The transport number of the silver ion ($t_{\ce{Ag+}}$).
    \item The liquid junction potential ($E_{lj}$) established at the boundary.
\end{enumerate}
\tcblower
\textbf{Part (a): Transport Number of $\ce{NO3-}$ ($t_-$):}
Recall the relationship between the EMF with transference and without transference:
\begin{equation*}
    E_{\text{wt}} = 2 t_- E_{\text{wot}}
\end{equation*}
Given: $E_{\text{wt}} = 0.0632\text{ V}$ and $E_{\text{wot}} = 0.0591\text{ V}$.
\begin{equation*}
    t_- = \frac{E_{\text{wt}}}{2 E_{\text{wot}}} = \frac{0.0632}{2 \times 0.0591} = \frac{0.0632}{0.1182} \approx 0.5347 \quad (\approx 0.535)
\end{equation*}

\textbf{Part (b): Transport Number of $\ce{Ag+}$ ($t_+$):}
\begin{equation*}
    t_+ = 1 - t_- = 1 - 0.5347 = 0.4653 \quad (\approx 0.465)
\end{equation*}

\textbf{Part (c): Liquid Junction Potential ($E_{lj}$):}
\begin{equation*}
    E_{lj} = E_{\text{wt}} - E_{\text{wot}} = 0.0632\text{ V} - 0.0591\text{ V} = +0.0041\text{ V} = +4.1\text{ mV}
\end{equation*}
\textbf{Physical Confirmation:}
Since $t_- > t_+$, the nitrate anion moves faster than the silver cation. The faster diffusion of $\ce{NO3-}$ into the dilute compartment charges it negatively, creating a positive junction potential that accelerates $\ce{Ag+}$ and retards $\ce{NO3-}$.
\end{examplebox}
